\documentclass[12pt]{article}
\usepackage[a4paper,margin=1in]{geometry}
\usepackage[utf8]{inputenc}
\usepackage[T1]{fontenc}
\usepackage{amsmath,amssymb}
\usepackage{textcomp}
\usepackage{parskip}
\begin{document}
\begin{center}
{\Large \textbf{2018 STEP I}}
\end{center}

\subsection*{Page 1}
\begin{flushleft}
\bigskip\textbf{Section A: Pure Mathematics}\par
\bigskip\textbf{Question 1.} The line y = a2x and the curve y = x(b x)2, where 0 < a < b , intersect at the origin O\par
-\par
and at points P and Q. The x-coordinate of P is less than the x-coordinate of Q. Find the\par
coordinates of P and Q, and sketch the line and the curve on the same axes.\par
Show that the equation of the tangent to the curve at P is\par
\[y = a(3a 2b)x + 2a(b a) .\]
- -\par
This tangent meets the y-axis at R. The area of the region between the curve and the line\par
segment OP is denoted by S. Show that\par
\[S = (b a) (3a + b) .\]
12 -\par
\[The area of triangle OP R is denoted by T . Show that S > 1 T .\]
log c\par
\[2 If x = log c , express c in terms of b and x and prove that a = log c .\]
b log b b\par
a\par
\textbf{(i)} Given that \pi2 < 10 , prove that\par
1 1\par
\[+ > 2 .\]
\[log \pi log \pi\]
2 5\par
\[\pi 1\]
\textbf{(ii)} Given that log > and that e2 < 8, prove that ln \pi > 17 .\par
2 e 5 15\par
\textbf{(iii)} Given that e3 > 20, \pi2 < 10 and log 2 > 3 , prove that ln \pi < 15 .\par
10 10 13\par
\bigskip\textbf{Question 3.} The points R and S have coordinates ( a , 0) and (2a , 0), respectively, where a > 0 . The\par
-\par
\[point P has coordinates (x , y) where y > 0 and x < 2a . Let ∠P RS = α and ∠P SR = β .\]
\textbf{(i)} Show that, if β = 2α , then P lies on the curve y2 = 3(x2 a2) .\par
-\par
\textbf{(ii)} Find the possible relationships between α and β when 0 < α < \pi and P lies on the\par
\[curve y^2 = 3(x^2 a^2) .\]
-\par
\end{flushleft}
\subsection*{Page 2}
\begin{flushleft}
\bigskip\textbf{Question 4.} The function f is defined by\par
1 2 2\par
\[f(x) = 1 (ln x) (x > 0, x = 1) .\]
\[x ln x - \leq\]
\[\leq \geq\]
\[Show that, when (ln x)2 = 1 , both f(x) = 0 and f \leq (x) = 0 .\]
The function F is defined by\par
x\par
\[f(t) dt for 0 < x < 1 ,\]
⎧(1/e\par
\[F(x) = ⎪ ⎪\]
⎪\par
⎪\par
⎨ x\par
\[f(t) dt for x > 1 .\]
⎪(e\par
⎪\par
⎪\par
⎪\par
⎩\par
\textbf{(i)} Find F(x) explicitly and hence show that F(x\par
-1)\par
\[= F(x) .\]
\textbf{(ii)} Sketch the curve with equation y = F(x) .\par
5 (i) Write down the most general polynomial of degree 4 that leaves a remainder of 1 when\par
divided by any of x 1 , x 2 , x 3 or x 4 .\par
- - - -\par
\textbf{(ii)} The polynomial P(x) has degree N , where N \geq 1 , and satisfies\par
\[P(1) = P(2) = = P(N ) = 1 .\]
\cdots\par
\[Show that P(N + 1) = 1 .\]
\[\leq\]
\[Given that P(N + 1) = 2 , find P(N + r) where r is a positive integer. Find a positive\]
\[integer r, independent of N , such that P(N + r) = N + r .\]
\textbf{(iii)} The polynomial S(x) has degree 4. It has integer coefficients and the coefficient of x4\par
is 1. It satisfies\par
\[S(a) = S(b) = S(c) = S(d) = 2001 ,\]
where a, b, c and d are distinct (not necessarily positive) integers.\par
\[(a) Show that there is no integer e such that S(e) = 2018 .\]
(b) Find the number of ways the (distinct) integers a, b, c and d can be chosen such\par
\[that S(0) = 2017 and a < b < c < d .\]
\end{flushleft}
\subsection*{Page 3}
\begin{flushleft}
\bigskip\textbf{Question 6.} Use the identity\par
\[2 sin P sin Q = cos(Q P ) cos(Q + P )\]
- -\par
to show that\par
\[2 sin θ sin θ + sin 3θ + + sin(2n 1)θ = 1 cos 2nθ .\]
\cdots - -\par
\[\leq \geq\]
\textbf{(i)} Let A be the approximation to the area under the curve y = sin x from x = 0 to\par
n\par
\[\pi\]
\[x = \pi, using n rectangular strips each of width , such that the midpoint of the top\]
n\par
of each strip lies on the curve. Show that\par
\[\pi \pi\]
\[A sin = .\]
n\par
2n n\par
) [\par
\textbf{(ii)} Let B be the approximation to the area under the curve y = sin x from x = 0 to\par
n\par
\[\pi\]
\[x = \pi , using the trapezium rule with n strips each of width . Show that\]
n\par
\[\pi \pi \pi\]
\[B sin = cos .\]
n\par
2n n 2n\par
) [ ) [\par
\textbf{(iii)} Show that\par
\[(A n + B n ) = B 2n ,\]
and that\par
\[A n B 2n = A 2n .\]
\[7 (i) In the cubic equation x^3 3pqx + pq(p + q) = 0 , where p and q are distinct real\]
-\par
numbers, use the substitution\par
pz + q\par
\[x =\]
z + 1\par
\[to show that the equation reduces to az3 + b = 0 , where a and b are to be expressed\]
in terms of p and q.\par
\textbf{(ii)} Show further that the equation x3 3cx + d = 0 , where c and d are non-zero real\par
\[numbers, can be written in the form - x^3 3pqx + pq(p + q) = 0 , where p and q are\]
\[distinct real numbers, provided d^2 > 4c3 . -\]
\textbf{(iii)} Find the real root of the cubic equation x3 + 6x 2 = 0 .\par
-\par
\textbf{(iv)} Find the roots of the equation x3 3p2x + 2p3 = 0 , and hence show how the equation\par
\[x^3 3cx + d = 0 can be solved in - the case d^2 = 4c3 .\]
-\par
\end{flushleft}
\subsection*{Page 4}
\begin{flushleft}
\bigskip\textbf{Question 8.} The functions s and c satisfy s(0) = 0 , c(0) = 1 and\par
\[\leq 2\]
\[s (x) = c(x) ,\]
\[\leq 2\]
\[c (x) = s(x) .\]
-\par
You may assume that s and c are uniquely defined by these conditions.\par
\textbf{(i)} Show that s(x)3 + c(x)3 is constant, and deduce that s(x)3 + c(x)3 = 1 .\par
\textbf{(ii)} Show that\par
d\par
\[s(x)c(x) = 2c(x) 1\]
dx -\par
) [\par
d s(x)\par
and find (and simplify) an expression in terms of c(x) for .\par
dx ] c(x) \{\par
\textbf{(iii)} Find the integrals\par
2 5\par
s(x) dx and s(x) dx .\par
( (\par
\textbf{(iv)} Given that s has an inverse function, s\par
-1,\par
\[use the substitution u = s(x) to show that\]
-1\par
\[du = s (u) + constant.\]
\[( (1 u^3) 3\]
-\par
\textbf{(v)} Find, in terms of u, the integrals\par
1 1\par
du and (1 u 3 ) 3 du .\par
\[( (1 u^3) 4 3 ( -\]
-\par
\end{flushleft}
\subsection*{Page 5}
\begin{flushleft}
\bigskip\textbf{Section B: Mechanics}\par
\bigskip\textbf{Question 9.} A straight road leading to my house consists of two sections. The first section is inclined\par
downwards at a constant angle α to the horizontal and ends in traffic lights; the second\par
section is inclined upwards at an angle β to the horizontal and ends at my house. The\par
distance between the traffic lights and my house is d.\par
I have a go-kart which I start from rest, pointing downhill, a distance x from the traffic lights\par
on the downward-sloping section. The go-kart is not powered in any way, all resistance forces\par
are negligible, and there is no sudden change of speed as I pass the traffic lights. Given that\par
\[I reach my house, show that x sin α \geq d sin β .\]
Let T be the total time taken to reach my house. Show that\par
g sin α 2\par
\[T = (1 + k)\sqrt\{x\} k2x kd ,\]
] 2 \{ - -\par
\}\par
sin α\par
\[where k = .\]
sin β\par
Hence determine, in terms of d and k, the value of x which minimises T . [You need not\par
justify the fact that the stationary value is a minimum.]\par
\bigskip\textbf{Question 10.} A train is made up of two engines, each of mass M , and n carriages, each of mass m. One of\par
the engines is at the front of the train, and the other is coupled between the kth and (k +1)th\par
carriages. When the train is accelerating along a straight, horizontal track, the resistance to\par
the motion of each carriage is R and the driving force on each engine is D, where 2D > nR .\par
The tension in the coupling between the engine at the front and the first carriage is T .\par
\textbf{(i)} Show that\par
n(mD + M R)\par
\[T = .\]
nm + 2M\par
\textbf{(ii)} Show that T is greater than the tension in any other coupling provided that k > 1 n .\par
\textbf{(iii)} Show also that, if k > 1 n , then at least one of the couplings is in compression (that\par
is, there is a negative tension in the coupling).\par
\end{flushleft}
\subsection*{Page 6}
\begin{flushleft}
\bigskip\textbf{Question 11.} The point O lies on a rough plane that is inclined at an angle α to the horizontal, where\par
\textasciicircum{}\circ\par
\[α < 45 . The point A lies on the plane a distance d from O up the line L of greatest slope\]
through O. The point B, which is not on the rough plane, lies in the same vertical plane as\par
O and A, and AB is horizontal. The distance from O to B is d.\par
A particle P of mass m rests on L between O and A. One end of a light inelastic string is\par
attached to P . The string passes over a smooth light pulley fixed at B and its other end is\par
attached to a freely hanging particle of mass λm.\par
\textbf{(i)} Show that the acute angle, θ, between the string and the line L satisfies α ( θ ( 2α .\par
\textbf{(ii)} Given that P can rest in equilibrium at every point on L between O and A, show that\par
2λ sin α ( 1 .\par
\textbf{(iii)} The coefficient of friction between P and the plane is μ, and the acute angle β is given\par
\[by μ = tan β . Show that if β \geq 2α , then a necessary condition for equilibrium to be\]
possible for every position of P on L between O and A is\par
sin(β α)\par
λ ( - .\par
cos(β 2α)\par
-\par
Obtain the corresponding result if α ( β ( 2α .\par
\end{flushleft}
\subsection*{Page 7}
\begin{flushleft}
\bigskip\textbf{Section C: Probability and Statistics}\par
\bigskip\textbf{Question 12.} A bag contains three coins. The probabilities of their showing heads when tossed are p 1, p 2\par
and p 3.\par
\textbf{(i)} A coin is taken at random from the bag and tossed. What is the probability that it\par
shows a head?\par
\textbf{(ii)} A coin is taken at random from the bag (containing three coins) and tossed; the coin\par
is returned to the bag and again a coin is taken at random from the bag and tossed.\par
Let N 1 be the random variable whose value is the number of heads shown on the two\par
tosses. Find the expectation of N 1 in terms of p, where p = 3 1 (p 1 + p 2 + p 3) , and show\par
\[that Var (N 1) = 2p(1 p) .\]
-\par
\textbf{(iii)} Two of the coins are taken at random from the bag (containing three coins) and tossed.\par
Let N 2 be the random variable whose value is the number of heads showing on the two\par
coins. Find E(N 2) and Var (N 2).\par
\textbf{(iv)} Show that Var (N 2) ( Var (N 1), with equality if and only if p 1 = p 2 = p 3 .\par
\bigskip\textbf{Question 13.} A multiple-choice test consists of five questions. For each question, n answers are given\par
(n \geq 2) only one of which is correct and candidates either attempt the question by choosing\par
one of the n given answers or do not attempt it.\par
For each question attempted, candidates receive two marks for the correct answer and lose\par
one mark for an incorrect answer. No marks are gained or lost for questions that are not\par
attempted. The pass mark is five.\par
Candidates A, B and C don't understand any of the questions so, for any question which\par
they attempt, they each choose one of the n given answers at random, independently of their\par
choices for any other question.\par
\textbf{(i)} Candidate A chooses in advance to attempt exactly k of the five questions, where\par
\[k = 0, 1, 2, 3, 4 or 5. Show that, in order to have the greatest probability of passing\]
\[the test, she should choose k = 4 .\]
\textbf{(ii)} Candidate B chooses at random the number of questions he will attempt, the six\par
possibilities being equally likely. Given that Candidate B passed the test find, in\par
terms of n, the probability that he attempted exactly four questions.\par
\textbf{(iii)} For each of the five questions Candidate C decides whether to attempt the question\par
by tossing a biased coin. The coin has a probability of n of showing a head, and\par
n+1\par
she attempts the question if it shows a head. Find the probability, in terms of n, that\par
Candidate C passes the test.\par
\end{flushleft}
\end{document}
